\section{Ergebnisse, Diskussion \& Fazit}\label{sec:ergebnisse_diskussion_fazit}

\subsection{Ergebnisse \& Evaluation}

\subsubsection{Was funktioniert}

Die Implementierung zeigt mehrere erfolgreiche Aspekte.
Der \textbf{Agent vs Human Modus} ermöglicht eine erfolgreiche Integration von trainierten Strategien aus der Bachelorarbeit.
Strategien werden korrekt geladen und der Agent spielt automatisch basierend auf der trainierten Strategie.
Die zufällige Zuweisung von Player-IDs sorgt dafür, dass der menschliche Spieler nicht immer derselbe Spieler ist, was die Spielerfahrung verbessert.
Die Integration funktioniert nahtlos, da die Strategien direkt aus der Bachelorarbeit wiederverwendet werden können, ohne dass eine Konvertierung erforderlich ist.

Der \textbf{Human vs Human Modus} funktioniert zuverlässig, wobei der WebSocket-Server Echtzeit-Multiplayer-Spiele ermöglicht.
Zwei Clients können sich verbinden und ohne merkbare Latenz gegeneinander spielen.
Die Push-basierte Kommunikation funktioniert wie erwartet, und State-Updates werden sofort an beide Clients gesendet.
Die automatische Wiederverbindung bei Verbindungsfehlern sorgt für eine robuste Verbindung.

Die \textbf{Wiederverwendung} der Game-Logik aus der Bachelorarbeit funktioniert ohne Änderungen, was die modulare Architektur bestätigt.
Die Game-Logik wurde erfolgreich in eine separate Klasse (\texttt{PokerGameLogic}) gekapselt, die sowohl von HTTP- als auch von WebSocket-Servern verwendet werden kann.
Dies ermöglicht Code-Wiederverwendung und vereinfacht die Wartung erheblich.

Die \textbf{Tests} bieten umfassende Test-Abdeckung für Server und Game-Logik, wobei alle automatisierten Tests bestehen.
Die szenarien-basierten Tests sind verständlich und geben guten Aufschluss über die Funktionalität des Systems.
Die Tests können ohne GUI-Fenster ausgeführt werden, was die Integration in CI/CD-Pipelines ermöglicht.

Die \textbf{Thread-Safety} wurde erfolgreich implementiert, wobei keine Race Conditions bei Multiplayer-Server beobachtet wurden.
\texttt{threading.Lock} verhindert erfolgreich gleichzeitige Zugriffe auf kritische Datenstrukturen.
Die Lock-Hold-Zeit ist kurz, da kritische Abschnitte auf das Minimum beschränkt wurden, was den Performance-Overhead minimiert.

\textbf{Sicherheitsfeatures} wie Input-Validierung und Rate Limiting funktionieren wie erwartet.
Input-Validierung verhindert erfolgreich XSS-Angriffe durch HTML-Escaping, und Rate Limiting schützt vor Missbrauch durch zu viele Requests.

\subsubsection{Einschränkungen}

Es gibt jedoch auch Einschränkungen.
Der Multiplayer-Modus unterstützt aktuell nur Heads-Up (2 Spieler).
Eine Erweiterung auf mehr Spieler würde Änderungen an der Game-Logik erfordern, da die aktuelle Implementierung für zwei Spieler optimiert ist.

Spielstände werden nicht gespeichert, sodass nach einem Neustart des Servers alle Spielzustände verloren gehen.
Dies ist eine bewusste Design-Entscheidung für diese Version, da Persistenz zusätzliche Komplexität erfordern würde.

Es gibt keine Replay-Funktionalität, sodass Spiele nicht nachträglich angeschaut werden können.
Dies wäre eine nützliche Funktion für Analyse und Lernen, würde jedoch zusätzliche Implementierung erfordern.

Große Strategie-Dateien benötigen Zeit zum Laden, was besonders bei Strategien mit vielen Iterationen spürbar ist.
Die Kompression hilft zwar, die Dateigröße zu reduzieren, aber das Laden und Dekomprimieren großer Dateien kann einige Sekunden dauern.
Dies ist jedoch akzeptabel, da Strategien nur einmal beim Start geladen werden.

Der Server hat keine Authentifizierung implementiert, sodass jeder sich verbinden kann, solange der Server nicht voll ist.
Dies ist für lokale Tests ausreichend, würde jedoch für Produktions-Einsatz zusätzliche Sicherheitsfeatures erfordern.

\subsubsection{Performance-Beobachtungen}

WebSocket zeigt deutlich niedrigere Latenz als HTTP-Polling.
Bei HTTP-Polling mit 100ms Intervall ist die Latenz mindestens 100ms, da Updates nur beim nächsten Polling-Intervall empfangen werden.
WebSocket hingegen ermöglicht sofortige Übertragung von Updates, was die Latenz praktisch eliminiert.
Dies verbessert die Spielerfahrung erheblich, da State-Updates sofort übertragen werden.

Komprimierte \texttt{.pkl.gz} Dateien reduzieren den Speicherbedarf erheblich.
Eine unkomprimierte Strategie-Datei kann mehrere hundert MB groß sein, während die komprimierte Version oft unter 50 MB liegt.
Die Ladezeit bleibt akzeptabel (wenige Sekunden), da Strategien nur einmal beim Start geladen werden.

\texttt{threading.Lock} verhindert Race Conditions effektiv und hat minimalen Performance-Overhead.
Die Lock-Hold-Zeit ist kurz, da kritische Abschnitte auf das Minimum beschränkt wurden.
Selbst bei vielen gleichzeitigen Requests wurde keine merkbare Verlangsamung beobachtet.

Die GUI bleibt auch während Netzwerk-Operationen responsiv, da \texttt{QThread} für WebSocket-Kommunikation verwendet wird und die GUI-Thread nicht blockiert wird.
Dies ermöglicht eine flüssige Benutzererfahrung, auch wenn Netzwerk-Operationen im Hintergrund stattfinden.

\subsubsection{Test-Ergebnisse}

Alle automatisierten Tests bestehen zuverlässig.
Während der Testphase wurden keine Bugs in der Server-Logik gefunden.
Die Input-Validierung und Rate Limiting funktionieren wie erwartet.
Die HTML-Escaping-Funktionalität verhindert erfolgreich XSS-Angriffe.
Tests mit bösartigen Eingaben (z.B. \texttt{<script>alert('XSS')</script>}) werden korrekt escaped und stellen keine Sicherheitsrisiko dar.

Rate Limiting funktioniert korrekt.
Clients, die zu viele Requests senden, werden abgelehnt, während normale Nutzung nicht beeinträchtigt wird.
Dies schützt den Server vor Missbrauch und stellt sicher, dass die Ressourcen fair verteilt werden.

Die Tests decken die wichtigsten Szenarien ab: Client-Registrierung, State-Updates, Action-Verarbeitung, Reset-Funktionalität, Fehlerbehandlung und Sicherheitsfeatures.
Die szenarien-basierten Tests sind verständlich und geben guten Aufschluss über die Funktionalität des Systems.

\subsection{Diskussion}

\subsubsection{Technische Bewertung}

Die modulare Architektur mit klarer Trennung von GUI, Client, Server und Logik erleichtert die Wartung erheblich.
Änderungen in einer Schicht beeinflussen andere Schichten minimal.
Beispielsweise kann die Game-Logik geändert werden, ohne dass die GUI oder die Client-Implementierungen angepasst werden müssen.
Dies ermöglicht unabhängige Entwicklung und Wartung der verschiedenen Komponenten.

Wiederverwendbare Komponenten folgen dem DRY-Prinzip, was Code-Duplikation vermeidet und die Wartung erleichtert.
Beide Spielmodi nutzen die gleichen UI-Komponenten, sodass Änderungen an diesen Komponenten automatisch für beide Modi gelten.
Dies reduziert den Wartungsaufwand erheblich.

Szenarien-basierte Tests sind verständlicher als viele Edge-Case-Tests und geben besseren Aufschluss über das tatsächliche Verhalten.
Die Tests decken typische Anwendungsfälle ab, wie sie in der Praxis auftreten würden, was das Debugging erleichtert.
Die Verwendung von Fixtures ermöglicht wiederverwendbare Test-Setups, was Code-Duplikation vermeidet.

Thread-Safety, Input-Validierung und Rate Limiting erhöhen die Robustheit des Systems und reduzieren die Wahrscheinlichkeit von Bugs.
Die Verwendung von Kontext-Managern für Thread-Safety stellt sicher, dass Locks auch bei Exceptions korrekt freigegeben werden.

\subsubsection{Design-Entscheidungen}

Die wichtigsten Design-Entscheidungen wurden bewusst getroffen.
PyQt6 wurde als mature Framework mit Signal-Slot-Pattern gewählt, was lose Kopplung zwischen Komponenten ermöglicht.
Dies erleichtert die Wartung und Erweiterung erheblich, da Komponenten unabhängig voneinander entwickelt werden können.

WebSocket wurde für Push-basierte Kommunikation für Echtzeit-Multiplayer gewählt, da es deutlich effizienter ist als HTTP-Polling.
Die niedrige Latenz verbessert die Spielerfahrung erheblich, und die Push-basierte Kommunikation reduziert den Netzwerk-Traffic.

pytest wurde gegenüber unittest gewählt, da es bessere Test-Organisation bietet.
Fixtures ermöglichen wiederverwendbare Test-Setups, und die Syntax ist lesbarer als unittest, was die Wartung der Tests erleichtert.

\texttt{threading.Lock} wurde als einfache, effektive Thread-Safety-Lösung gewählt.
Es ist einfach zu verwenden, hat minimalen Performance-Overhead und ist für die Anforderungen dieses Projekts ausreichend.
Komplexere Ansätze wie Actor-Model oder Queue-basierte Patterns wären für diesen Anwendungsfall Overkill gewesen.

\subsubsection{Erkenntnisse und Advanced Python Konzepte}

Die Implementierung demonstriert mehrere wichtige Advanced Python Konzepte.
PyQt6 Signals/Slots ermöglichen lose Kopplung zwischen Komponenten und implementieren das Observer-Pattern.
Komponenten können Signale emittieren, ohne zu wissen, wer diese empfängt, was die Wartbarkeit und Erweiterbarkeit erheblich verbessert.

Wiederverwendbare GUI-Komponenten folgen dem DRY-Prinzip für verschiedene Spielmodi.
Beide Modi nutzen die gleichen UI-Komponenten, was Code-Duplikation vermeidet und die Wartung erleichtert.

Asyncio + QThread Integration ermöglicht eine Bridge zwischen async und Qt Event Loop, was Push-basierte Kommunikation statt Polling ermöglicht.
Die Lösung mit \texttt{WebSocketWorker} als \texttt{QThread} mit eigener asyncio-Event-Loop zeigt, wie verschiedene Event-Loop-Systeme integriert werden können.
Die Signal-Bridge überträgt Async-Events via PyQt6 Signals an den GUI-Thread, wodurch eine saubere Trennung zwischen asynchroner Netzwerk-Kommunikation und GUI-Updates erreicht wird.

Thread-Safety mit \texttt{threading.Lock} und Kontext-Managern stellt sicher, dass kritische Operationen Thread-Safe sind.
Die Verwendung von Kontext-Managern (\texttt{with self.lock:}) stellt sicher, dass Locks auch bei Exceptions korrekt freigegeben werden.
Dies ist essentiell für Multiplayer-Server, bei denen mehrere Clients gleichzeitig auf gemeinsame Datenstrukturen zugreifen.

Decorators ermöglichen Rate Limiting als wiederverwendbare Funktionalität.
Die Implementierung als Decorator folgt dem Decorator-Pattern und macht Rate Limiting einfach anwendbar.
Dies demonstriert die Mächtigkeit von Python-Decorators für Cross-Cutting Concerns wie Sicherheit und Performance.

Die erfolgreiche Wiederverwendung der Game-Logik aus der Bachelorarbeit zeigt, dass modulare Architekturen Wiederverwendung ermöglichen.
Die Game-Logik wurde erfolgreich in eine separate Klasse gekapselt, die sowohl von HTTP- als auch von WebSocket-Servern verwendet werden kann.
Dies bestätigt die Wichtigkeit einer guten Architektur für Code-Wiederverwendung.

\subsection{Fazit}

Die Implementierung einer Poker-GUI mit PyQt6 für Agent vs Human und Multiplayer-Modi wurde erfolgreich abgeschlossen.
Beide Spielmodi funktionieren zuverlässig und bieten eine gute Spielerfahrung.
Die modulare Architektur ermöglicht Wiederverwendung von Code aus der Bachelorarbeit und erleichtert die Wartung erheblich.

Die Verwendung von WebSocket für Echtzeit-Kommunikation reduziert die Latenz erheblich im Vergleich zu HTTP-Polling.
Thread-Safety wurde erfolgreich implementiert, und alle automatisierten Tests bestehen zuverlässig.
Die Implementierung demonstriert mehrere wichtige Advanced Python Konzepte wie Signals/Slots, asyncio-Integration, Thread-Safety und Decorators.

\subsubsection{Ausblick}

Mögliche Erweiterungen für zukünftige Versionen umfassen Unterstützung für mehr Spieler, einen Turnier-Modus mit mehreren Runden, Replay-Funktionalität zur Aufzeichnung und Wiedergabe von Spielen, Performance-Optimierungen wie Strategie-Caching und Lazy Loading, erweiterte Sicherheitsfeatures wie Authentifizierung und Verschlüsselung sowie UI-Verbesserungen wie bessere Animationen und Themes.

Die modulare Architektur ermöglicht diese Erweiterungen ohne größere Umstrukturierungen.
Die klare Trennung zwischen GUI, Client, Server und Logik ermöglicht es, einzelne Komponenten zu erweitern, ohne andere Komponenten zu beeinflussen.
Dies zeigt, wie wichtig eine gute Architektur für die langfristige Wartbarkeit und Erweiterbarkeit eines Systems ist.
